%!TEX root = ../Quang Anh_FinalProject.tex
% --- Cài đặt ngôn ngữ và font chữ ---
\usepackage{iftex}
\ifPDFTeX
  \usepackage[utf8]{vietnam}
  \usepackage{tipa}
  \usepackage{newunicodechar}
  \newunicodechar{θ}{\texttheta}
  \newunicodechar{ð}{\texteth}
  \newunicodechar{ʃ}{\textesh}
  \newunicodechar{ʒ}{\textyogh}
  \newunicodechar{ŋ}{\texteng}
  \newunicodechar{ɡ}{\textg}
  \newunicodechar{ʔ}{\textglotstop}
  \newunicodechar{ɪ}{\textsci}
  \newunicodechar{ː}{\textlengthmark}
  \newunicodechar{æ}{\textae}
  \newunicodechar{ʌ}{\textturnv}
  \newunicodechar{ɒ}{\textturnscripta}
  \newunicodechar{ɔ}{\textopeno}
  \newunicodechar{ə}{\textschwa}
  \newunicodechar{ʊ}{\textupsilon}
  \newunicodechar{ɛ}{\textepsilon}
  \newunicodechar{ɜ}{\textrevepsilon}
  \newunicodechar{ˈ}{\textprimstress}
  \newunicodechar{ˌ}{\textsecstress}
  \newcommand{\ipa}[1]{\textnormal{#1}}
\else
  \usepackage{fontspec}
  \usepackage[vietnamese]{babel}
  \setmainfont{lmroman10-regular.otf}[
    BoldFont = lmroman10-bold.otf,
    ItalicFont = lmroman10-italic.otf,
    BoldItalicFont = lmroman10-bolditalic.otf
  ]
  \IfFileExists{setup/fonts/Doulos_SIL_Regular.ttf}{%
    \newfontfamily{\doulossil}{Doulos_SIL_Regular.ttf}[
      Path = setup/fonts/,
      Scale = MatchLowercase
    ]%
  }{%
    \newcommand{\doulossil}{}%
  }
  \newcommand{\ipa}[1]{{\doulossil #1}}
\fi
\usepackage{anyfontsize}
\usepackage{microtype}

% --- Định dạng trang ---
% 1. Cấu hình 1 mặt (oneside): Lề trái 3.5cm (đóng gáy), lề phải 2.0cm, trên 3.5cm, dưới 3.0cm
\usepackage[left=3.5cm,right=2cm,top=3.5cm,bottom=3cm,headheight=15pt]{geometry}

% 2. Cấu hình 2 mặt (twoside): Mở comment dòng dưới và đóng comment dòng trên khi in 2 mặt
%\usepackage[inner=3.5cm,outer=2cm,top=3.5cm,bottom=3cm,headheight=15pt,twoside]{geometry}
\usepackage{indentfirst} % Thụt lề đầu dòng

% --- Toán học ---
\usepackage{amsmath, amssymb, latexsym, amscd, amsxtra, amsthm}
\usepackage{mathtools}
\usepackage{pb-diagram}
\usepackage{aligned-overset}
%\usepackage{commath}

% --- Đồ họa và Hình ảnh ---
\usepackage{graphicx}
\usepackage{float}
\usepackage{wrapfig}
\usepackage{picinpar, floatflt}
\usepackage{caption}
\usepackage{pdfpages} % Chèn file PDF
\usepackage{tikz, pgf}
\usetikzlibrary{shapes, patterns, arrows.meta, calc}

% --- Bảng biểu ---
\usepackage{multirow}
\usepackage{longtable}
\usepackage{tabularx}
\usepackage{booktabs}
\newcolumntype{L}{>{\raggedright\arraybackslash}X}

% --- Tiện ích khác ---
\usepackage{enumitem}
\usepackage{fancybox}
\usepackage{comment} % Gói lệnh comment
\usepackage[tight,vietnam]{minitoc}
\usepackage{hyperref}

% --- Code Snippet ---
\usepackage{listings}
\usepackage{xcolor}
\definecolor{codegreen}{rgb}{0,0.6,0}
\definecolor{codegray}{rgb}{0.5,0.5,0.5}
\definecolor{codepurple}{rgb}{0.58,0,0.82}
\definecolor{backcolour}{rgb}{0.95,0.95,0.92}
\lstdefinestyle{mystyle}{
	backgroundcolor=\color{backcolour},   
	commentstyle=\color{codegreen},
	keywordstyle=\color{magenta},
	numberstyle=\tiny\color{codegray},
	stringstyle=\color{codepurple},
	basicstyle=\ttfamily\normalsize,
	breakatwhitespace=false,         
	breaklines=true,                 
	captionpos=b,                    
	keepspaces=true,                 
	numbers=left,                    
	numbersep=5pt,                  
	showspaces=false,                
	showstringspaces=false,
	showtabs=false,                  
	tabsize=2
}
\lstset{style=mystyle}

% --- Cấu hình Mục lục từng chương (TitleTOC) ---
\usepackage{titletoc}
\titlecontents{chapter}
  [0pt]
  {\bfseries}
  {\chaptername\ \thecontentslabel\enskip}
  {}
  {\hfill\contentspage}

% --- Quản lý Tài liệu tham khảo (biblatex) ---
\usepackage{csquotes}
\usepackage[backend=biber, style=numeric, sorting=nyt, defernumbers=true]{biblatex}
\addbibresource{backmatter/references.bib}
